% Threats to validity: last subsection of the evaluation (\input inside sections/evaluation).

\subsection{Threats to Validity}
\label{sec:threats}

\textbf{Construct Validity:}
We count the held-out instances on which a rule list is certified, which depends on the thresholds.
A tight threshold can hide differences between methods, as on Pac-Man, where the threshold lies just above the optimum.
The thresholds are the same for every method, taken from prior work for the grid world and set relative to each instance's optimum otherwise\todo{confirm they were fixed before the runs}, and on Pac-Man we also report the worst-case crash probability at each horizon.
A certificate is only as good as the computed values, so every reported value is recomputed with interval iteration, whose error is bounded\todo{some checks fall back to Gauss--Seidel at precision $10^{-12}$ where interval iteration does not converge; count them}.
Finally, the evaluation certifies a library seeded before deployment on held-out instances, without the library growth of \autoref{sec:approach}, so it measures how often the system can reuse a rule list rather than the full loop.

\textbf{Internal Validity:}
We designed each domain's vocabulary, which could favor our rule lists.
To mitigate this, the decision trees, reinforcement learning and the search over rule lists receive the same vocabulary, and we also report grid worlds without the shortest-path feature \texttt{closer}.
The configuration of every method, ours and the baselines', is chosen by the same rule on validation instances, and the results we report are on these instances, while the test instances stay sealed\todo{state whether the camera-ready reports the sealed test sets}.
Selection fixes the order among the actions of \texttt{any} rules on the training instances, so the model checker contributes this choice, and we therefore also report certification with the order left open.
LLM output varies between runs, and the hosted models take no seed, so we repeat each configuration three times and report the mean and the range.
We reimplemented 1--2--3--Go! and ran it both on its paper's training instances and on ours\todo{check the reimplementation against the published values}.

\textbf{External Validity:}
Our three domains vary the size, layout and horizon of their instances, and they combine reachability, ordering and safety requirements.
Pac-Man and Mars Exploration Rovers are public benchmarks that an LLM may have seen during training, whereas we generated the grid-world families for this paper.
We evaluate four LLMs\todo{name the providers once the runs are final} to reduce the dependence on one model.
Our certificates hold on the instances we checked, up to the sizes PRISM can build, and we make no claim about instances beyond them.
